\chapter{Ранга-один конденсат когерентности стрелок}
\label{ch:v7-edge-coherence-rank-one-condensate-gate}

\section{Синтез проверенной формульной базы}

Предыдущий гейт локализовал недостающий объект: две разрешённые стрелки
должны не складываться независимо, а образовать ранга-один когерентную
ковариацию. Проект уже проверил три близких механизма. Энтропия предпочитает
полный ранг, плоская мера частичных изометрий --- ранг два, а норма внешнего
квадрата обращается в нуль точно на ранге не выше одного. Поэтому новый тест
не повторяет энтропийный или мерный селектор.

Пусть три общих соседних канала собраны в поле
\begin{equation}
 B\in\operatorname{Hom}(\mathbb C_{\rm ch}^3,
 \mathbb C_{\rm copy}^2)\simeq M_{2\times3}(\mathbb C),
 \qquad C=BB^\dagger.
 \label{eq:v7-edge-coherence-field}
\end{equation}

\section{Прямой минус перекрёстный путь}

Определим антисимметризованную двухпутевую амплитуду
\begin{equation}
 W_{ia,jb}(B)=B_{ia}B_{jb}-B_{ib}B_{ja}.
 \label{eq:v7-edge-coherence-wedge-tensor}
\end{equation}
Тождество Коши--Бине даёт
\begin{equation}
 \|W(B)\|_F^2
 =4\|\Lambda^2B\|_F^2
 =4\det(BB^\dagger).
 \label{eq:v7-edge-coherence-wedge-identity}
\end{equation}
Коэффициент четыре возникает из полного суммирования прямого минус
перекрёстного тензора, как в ранговом селекторе Тома IV и его проверке в
Томе VI.

\section{Минимальный потенциал уровня конфигураций}

Число три уже задано количеством общих каналов. Рассмотрим
\begin{equation}
 \mathcal S_{\rm coh}(B)
 =\left(\Tr BB^\dagger-3\right)^2+\|W(B)\|_F^2.
 \label{eq:v7-edge-coherence-action}
\end{equation}
Оба слагаемых неотрицательны, поэтому
\begin{equation}
 \mathcal S_{\rm coh}(B)=0
 \quad\Longleftrightarrow\quad
 \Tr BB^\dagger=3,
 \qquad \rank B=1.
 \label{eq:v7-edge-coherence-zero-set}
\end{equation}
Все минимумы имеют сингулярную форму
\begin{equation}
 B_\star=\sqrt3\,u v^\dagger,
 \qquad u\in\mathbb C^2,\qquad v\in\mathbb C^3,
 \qquad \|u\|=\|v\|=1.
 \label{eq:v7-edge-coherence-vacuum}
\end{equation}
После удаления общей фазы вакуумное многообразие имеет вещественную
размерность
\begin{equation}
 \dim_{\mathbb R}\mathcal M_{\rm coh}
 =\dim S^3+\dim S^5-\dim U(1)=7.
 \label{eq:v7-edge-coherence-vacuum-dimension}
\end{equation}

\section{Эндогенный запуск и полный гессиан}

В нуле действие равно девяти, а его квадратичная часть равна
$-6\Tr BB^\dagger$. Поэтому полный вещественный гессиан на двенадцати
компонентах имеет спектр
\begin{equation}
 \Spec\Hess_0\mathcal S_{\rm coh}
 =\{-12\}^{\times12}.
 \label{eq:v7-edge-coherence-origin-hessian}
\end{equation}
Нулевой сектор стационарен и неустойчив без внешнего толчка.

В представителе
$B_0=\begin{psmallmatrix}\sqrt3&0&0\\0&0&0\end{psmallmatrix}$
машинный и аналитический расчёты дают
\begin{equation}
 \Spec\Hess_{B_0}\mathcal S_{\rm coh}
 =\{0\}^{\times7}\cup\{24\}^{\times5}.
 \label{eq:v7-edge-coherence-vacuum-hessian}
\end{equation}
Семь нулей точно касательны к $\mathcal M_{\rm coh}$; отрицательных
поперечных мод нет.

\section{Индуцированный проектор и модулярная высота}

Редуцированное копийное состояние автоматически чисто:
\begin{equation}
 R_{\rm copy}(B_\star)
 =\frac{B_\star B_\star^\dagger}{\Tr B_\star B_\star^\dagger}
 =uu^\dagger,
 \qquad R_{\rm copy}^2=R_{\rm copy}.
 \label{eq:v7-edge-coherence-copy-projector}
\end{equation}
Следовательно, высота
\begin{equation}
 h_{\rm copy}=2R_{\rm copy}-I_2,
 \qquad h_{\rm copy}^2=I_2
 \label{eq:v7-edge-coherence-height}
\end{equation}
возникает после конденсации, а не вставляется как фиксированная
$\sigma_x$. Модулярный расчёт предыдущего гейта переносится на любую
спонтанно выбранную ось $u$.

Правое редуцированное состояние также имеет ранг один:
\begin{equation}
 R_{\rm ch}(B_\star)
 =\frac{B_\star^\dagger B_\star}{\Tr B_\star^\dagger B_\star}
 =vv^\dagger.
 \label{eq:v7-edge-coherence-channel-projector}
\end{equation}
Поэтому потенциал одновременно выбирает одну копийную линию и одну
когерентную комбинацию трёх соседних каналов. Вторая ориентация пока не
сопоставлена наблюдаемому меню рёбер.

\section{Real-ветвь и граница результата}

Для вещественного $B$ множество минимумов равно
$(S^1\times S^2)/\mathbb Z_2$ и имеет размерность три. Real-условие уменьшает
число плоских направлений, но не выбирает конкретные $u$ и $v$.

Главный разрыв остаётся родительским. Радиальный квадрат может наследовать
замкнутый мост и число общих каналов, а внешний квадрат имеет точный
double-path прецедент. Но в текущей конечной геометрии ещё не построена одна
суперсвязность, чья единственная норма кривизны одновременно содержит оба
слагаемых~\eqref{eq:v7-edge-coherence-action} с указанной относительной
нормировкой. Простое добавление нормы миноров было бы новым ручным сектором.

\begin{theorem}[Условный конденсат стрелочной когерентности]
Функционал~\eqref{eq:v7-edge-coherence-action} имеет неустойчивый нулевой
сектор и точное многообразие поперечно устойчивых ненулевых минимумов ранга
один. Каждый минимум порождает чистый копийный проектор и модулярную высоту
без заранее выбранной оси. Это закрывает задачу на уровне потенциала, но не
выводит потенциал из одного конечного родителя.
\end{theorem}

\begin{nogocriterion}[Запрет ручной нормы миноров]
Нельзя считать сумму радиального квадрата и $\|W\|^2$ единым физическим
действием, пока прямой и перекрёстный пути, их относительный минус и общая
следовая метрика не получены из одной Real-градуированной архитектуры.
Нельзя также называть выбранные линии физическими поколениями или частицами
до проверки полного носителя и калибровочных токов.
\end{nogocriterion}

\begin{protocol}
\begin{enumerate}
 \item поле когерентности $2\times3$: \textbf{построено};
 \item точный внешний квадрат: \textbf{получен};
 \item нулевой сектор: \textbf{стационарен и имеет 12 отрицательных мод};
 \item ненулевой минимум: \textbf{ранг один, норма три};
 \item полный поперечный гессиан: \textbf{пять положительных мод};
 \item касательные нули: \textbf{семь, совпадают с размерностью орбиты};
 \item чистый копийный проектор: \textbf{возникает};
 \item фиксированная внешняя ось: \textbf{не вводится};
 \item один кривизностный родитель обоих слагаемых: \textbf{не построен};
 \item итог: \textbf{условно положительный потенциал};
 \item следующий гейт: \textbf{внешнеквадратный родитель когерентности и
 единая следовая нормировка}.
\end{enumerate}
\end{protocol}

Машинный сертификат сохранён в
\path{s2t_v7_edge_coherence_rank_one_condensate_gate_results.json}.

\begin{thebibliography}{9}
\bibitem{v7coherence-quiver}
M.~Harada, G.~Wilkin,
\emph{Morse theory of the moment map for representations of quivers},
Geometriae Dedicata 150 (2011), 307--353; arXiv:0807.4734.
\bibitem{v7coherence-krajewski}
T.~Krajewski,
\emph{Classification of finite spectral triples},
Journal of Geometry and Physics 28 (1998), 1--30;
arXiv:hep-th/9701081.
\end{thebibliography}